%=========== ΛΥΣΗ ΘΕΜΑΤΟΣ Γ.77 ===========
\begin{thema}{Γ}
Δίνεται $f^2(x) - 4f(x) = x^2 - 4$ (1), $f(0)=2$ και $f(x) \ge 2$.
\begin{erwthma}

\item Αναδιατάσσουμε την εξίσωση (1) συμπληρώνοντας το τέλειο τετράγωνο στο πρώτο μέλος. Προσθέτουμε το $4$ και στα δύο μέλη:
\begin{align*}
f^2(x) - 4f(x) + 4 &= x^2 - 4 + 4 \\
(f(x) - 2)^2 &= x^2
\end{align*}
Εξάγουμε την τετραγωνική ρίζα:
\[ |f(x) - 2| = \sqrt{x^2} \Rightarrow |f(x) - 2| = |x| \]
Από τα δεδομένα γνωρίζουμε ρητά ότι $f(x) \ge 2 \Rightarrow f(x) - 2 \ge 0$. Συνεπώς το απόλυτο στο αριστερό μέλος καταργείται με θετικό πρόσημο ανεξαρτήτως:
\[ f(x) - 2 = |x| \Rightarrow \mathbf{f(x) = |x| + 2} \]
(Η επαλήθευση είναι άμεση: $f(0) = |0|+2 = 2$).

\item Ζητείται να αποδείξουμε την παραγωγισιμότητα στο $x_0=0$.
Υπολογίζουμε τις πλευρικές παραγώγους μέσω του ορισμού (του ορίου του λόγου μεταβολής) στο 0:
- Για $x > 0$: $\lim_{x\to 0^+} \frac{f(x) - f(0)}{x - 0} = \lim_{x\to 0^+} \frac{x + 2 - 2}{x} = \lim_{x\to 0^+} \frac{x}{x} = 1$.
- Για $x < 0$: $\lim_{x\to 0^-} \frac{f(x) - f(0)}{x - 0} = \lim_{x\to 0^-} \frac{-x + 2 - 2}{x} = \lim_{x\to 0^-} \frac{-x}{x} = -1$.
Εφόσον η αριστερή παράγωγος ($-1$) διαφέρει από τη δεξιά ($+1$), η συνάρτηση \textbf{δεν είναι παραγωγίσιμη} στο $x_0=0$ (αποτελεί γωνιακό σημείο).

\begin{parat}{\linewidth}
\textbf{Σφάλμα Εκφώνησης:} Η διατύπωση παρασίτησε ζητώντας «Να αποδείξετε ότι ΕΙΝΑΙ παραγωγίσιμη», κάτι που, όπως αποδείξαμε ακλόνητα, είναι απόλυτα λανθασμένο μαθηματικά, μιας και η συνάρτηση τέμνεται ως απόλυτη τιμή.
\end{parat}

\item Ένα τυχαίο σημείο της καμπύλης έχει συντεταγμένες $M(x, f(x))$.
Το τετράγωνο της απόστασης του σημείου αυτού από την αρχή των αξόνων $O(0,0)$ δίνεται από τον τύπο $d^2(x) = x^2 + [f(x)]^2$.
Μας ζητείται να εντοπίσουμε ελάχιστο \textbf{αυστηρά για $x > 0$}.
Για $x > 0$, ισχύει $f(x) = x + 2$. Η συνάρτηση απόστασης (στο τετράγωνο) γίνεται:
\[ g(x) = d^2(x) = x^2 + (x+2)^2 = 2x^2 + 4x + 4 \]
Παραγωγίζουμε για να μελετήσουμε τη μονοτονία:
\[ g'(x) = 4x + 4 = 4(x+1) \]
Εφόσον δεσμευόμαστε αποκλειστικά σε $x > 0$, η παράγωγος είναι μονίμως θετική ($g'(x) > 4 > 0$).
Συνεπώς η ταχύτητα ανόδου (και η απόσταση) είναι \textbf{γνησίως αύξουσα} σε ολόκληρο το ανοιχτό $(0, +\infty)$.
Διόλου περίεργα, το ελάχιστο «προσπαθεί» να υπάρξει στο όριο $x \to 0$, όμως το 0 δεν περιλαμβάνεται! Άρα \textbf{δεν υφίσταται κανένα γεωμετρικό σημείο ελαχίστου} (παρά μόνο αν κλείσουμε το διάστημα στο $x \ge 0$, όπου το $O(0,2)$ είναι η κατώτατη οριακή προσέγγιση της).

\item Προϋποθέσεις του Θεωρήματος \eng{Rolle} για $f(x) = |x|+2$ στο $[-1, 1]$:
1. Η $f$ είναι συνεχής στο κλειστό $[-1, 1]$.
2. $f(-1) = |-1|+2 = 3$ και $f(1) = |1|+2 = 3$. Άρα $f(-1) = f(1)$.
3. Είναι η $f$ παραγωγίσιμη στο εσωτερικό $(-1, 1)$; \textbf{ΟΧΙ}. Στο ενδιάμεσο σημείο $x_0=0$ (όπως ήδη επιλύθηκε στο Ερώτημα 2), η συνάρτηση αστοχεί να παραχθεί (δημιουργεί γωνία).

Επομένως, \textbf{δεν ικανοποιούνται πλήρως οι προϋποθέσεις} και το Θεώρημα \eng{Rolle} \textbf{δεν δύναται να εφαρμοστεί}.
\end{erwthma}
\end{thema}
